\section{Explicit adelic base points and supported modifications}
\label{sec:adelic-geometry}

\status{AD1 constructs the actual diagonal base point and its last-row
residues. AD2--AD3 prove the local transition and supported-splice
arguments from supplied realizers. AD4 installs prescribed rational
frames without changing the carried lattice or the other decorations.
The initial realization of a self-dual decorated lattice remains a
separate geometric input, as do the identities for the particular
up/down neighbors. No rational strong-approximation assertion is used.}

\paragraph{Conventions.}
Let $E/F$ be the supplied CM extension, write $\bar x$ for its
involution, and put $\mathbb A=E\otimes_F\mathbb A_{F,f}$.
Use its canonical identification with $\mathbb A_{E,f}$, so an
adelic equality can be checked at every finite prime $P$ of $E$.
Write $\widehat{\mathcal O}=\prod_P\mathcal O_{E_P}$ for its
integral subring. Conjugation interchanges the $P$ and $\bar P$
components, using the induced isomorphism of their completions.
A matrix $g\in\mathrm{GL}_n(\mathbb A)$ is unitary when
$\bar g^{\mathsf T}g=1$. It carries the standard lattice to $N$
when $g\widehat{\mathcal O}^{n}=\widehat N$.
The full rational lattice is recovered from these local modules
by LP2--LP3 and the completion argument proved in AD4 below.

For a finite prime $P$, the characteristic function $e_P$ of its
component is an integral finite adele. Thus
\[
 e_P^2=e_P,\qquad e_Pe_Q=0\ (P\ne Q),\qquad
 \overline{e_P}=e_{\bar P}.                         \tag{AD.1}
\]
These are literal elements of the restricted product: only one
component differs from zero. They give finite componentwise
constructions inside the same adelic algebra. The source's
\texttt{idem} is the same idempotent, written as the idele with
component $2$ at $P$ and $1$ elsewhere, minus $1$.

\paragraph{AD1: a diagonal unitary base point with the prescribed residue.}
Let $0\ne\pi\in\mathcal O_E$, put $q=\bar\pi$, and assume
$(\pi,q)=\mathcal O_E$. The finite set $S$ of primes dividing
$\pi$ is disjoint from $\bar S$: otherwise one prime would
contain both terms in a Bezout identity for $\pi,q$.
Define
\[
 a_P=
 \begin{cases}
 q^{-1},&P\in S,\\
 \pi,&\bar P\in S,\\
 1,&P\notin S\cup\bar S.
 \end{cases}                                         \tag{AD.2}
\]
All displayed field elements are mapped to their indicated
completions. This is a finite modification of $1$, hence a finite
adele. At a prime in $S$, $q$ is a local unit; at a prime in
$\bar S$, $\pi$ is a local unit. Consequently both $a$ and
$a^{-1}$ belong to $\widehat{\mathcal O}$. At $P\in S$,
$\bar a_P=q$, while at $P\in\bar S$, $\bar a_P=\pi^{-1}$.
Checking the three cases proves
\[
                 a\bar a=1,\qquad a^{-1}=\bar a.      \tag{AD.3}
\]
In the idempotent notation, (AD.2) is also the finite formula
\[
 a=1-\sum_{P\in S}(e_P+\bar e_P)
          +\sum_{P\in S}e_Pq^{-1}
          +\sum_{P\in S}\bar e_P\pi.
\]
This proves the literal \texttt{exists\_unit}, including the
integrality of both inverses and $a_P^{-1}=q$ at each selected prime.

For $n>0$ put
\[
 k=\operatorname{diag}(1,\ldots,1,a),\qquad
 k^{-1}=\operatorname{diag}(1,\ldots,1,\bar a).        \tag{AD.4}
\]
Equation (AD.3) proves that $k$ is unitary. Both matrices are
integral, so $k\widehat{\mathcal O}^{n}=
\widehat{\mathcal O}^{n}$. It therefore realizes the actual
standard lattice, with empty type, auxiliary and twist decoration.
If $\pi q=T\ne0$ is a rational integer, its components are $1$
at every prime not dividing $T$. All $E$-components over an
$F$-prime $v\nmid T$ are therefore $1$, so the projection of
$k$ to $E\otimes_FF_v$ is exactly the identity. At a selected
prime $P\mid\pi$, the last row of $k^{-1}$ is
\[
              (0,\ldots,0,q),                       \tag{AD.5}
\]
already before reduction. Thus its entries have the precise
\texttt{HasRes} values required in \texttt{exists\_base} and
\texttt{exists\_kStd}. In dimension zero take the unique empty
identity matrix; the last-row assertions have no instances.
The actual source hypotheses $T=T_u$, $\pi=\pi_u$ give this
construction for every word and every dimension, simultaneously.

\paragraph{AD2: local lattice inclusion controls transition matrices.}
Suppose $a,b\in\mathrm{GL}_n(\mathbb A)$ carry the standard
lattice to full lattices $N,N'$. If $N'_P\subseteq N_P$, then
\[
                      a_P^{-1}b_P\in M_n(\mathcal O_{E_P}).
                                                               \tag{AD.6}
\]
Indeed every column $b_Pe_j$ belongs to $N'_P$ and hence to
$a_P\mathcal O_{E_P}^n$; its coordinates under $a_P^{-1}$
are integral. Here a localization inclusion is first extended
to the completion. The extension is justified by clearing a
denominator on each element of a finite generating set.
Equivalently, for $x\in N'$ choose $s\notin P$ with
$sx\in N$; $s$ is a local unit, so integrality of
$a_P^{-1}(sx)$ implies integrality of $a_P^{-1}x$.
Finite sums and integral scalar multiples give the completed
module statement. If $N'_P=N_P$, apply the same reasoning in
the opposite direction to conclude
\[
                      a_P^{-1}b_P\in\mathrm{GL}_n(\mathcal O_{E_P}).
                                                               \tag{AD.7}
\]
This is the source's \texttt{valued\_inv\_mul\_le}, with the
inverse assertion available when the lattices agree.

One residue cancellation will also be needed. If $s\in\mathcal O_E
\setminus P$, $z\in E_P$, and $sz$ has residue $\bar s c$,
then $z$ is integral and has residue $c$. The first assertion
uses that $s$ is a unit; the second cancels its nonzero residue
in $\mathcal O_E/P$. In the source's explicit \texttt{HasRes}
notation, choose $t\in\mathcal O_E$ and $r\in P$ with
$ts+r=1$. If $sz-b$ has positive valuation, then
\[
             z-tb=t(sz-b)+rz                         \tag{AD.8}
\]
has positive valuation too, and the residue of $tb$ is $c$.
This proves \texttt{hasRes\_of\_mul}. It extends a residue
identity on a global lattice to any vector in its localization,
by multiplying the vector by a denominator outside $P$.

\paragraph{AD3: confine a realized modification to one conjugate pair.}
Let $X,X_0$ be the source's decorated lattices with type, auxiliary
and frame-prime sets $S_T,S_G,Q$. They are pairwise disjoint,
and their union contains no conjugate pair. Choose $P\in S_T$.
Suppose $h$ realizes $X$ and $h_0$ realizes $X_0$, in the full
\texttt{IsPtOf} sense. Assume that their localized lattices agree
away from $P,\bar P$, their type functionals agree at every
other prime of $S_T$, and their auxiliary reductions and rational
frames agree. These are the actual hypotheses used in the splice.

Put $e=e_P+e_{\bar P}$. It is an idempotent fixed by conjugation.
Define, entrywise,
\[
 g=(1-e)h+eh_0,\qquad
 g^{-1}=(1-e)h^{-1}+eh_0^{-1},\qquad
 \varepsilon=h^{-1}g.                                \tag{AD.9}
\]
The first two matrices are inverse because $e(1-e)=0$.
Since $\bar e=e$ and $h,h_0$ are unitary, the conjugate
transpose of the first is the second. Thus $g$ and
$\varepsilon$ are unitary. Their components satisfy
\[
 g_w=\begin{cases}(h_0)_w,&w=P,\bar P,\\h_w,&\text{otherwise},
       \end{cases}
 \qquad \varepsilon_w=1\quad(w\ne P,\bar P).          \tag{AD.10}
\]
There is no restriction on the size of the changed components;
only finitely many are changed.

We check all five fields of \texttt{IsPtOf} for $g,X_0$.
Unitarity was proved above. For the lattice field,
$g_w^{-1}(h_0)_w$ and its inverse are identity at the two
changed primes and are integral elsewhere by (AD.7).
Therefore $g^{-1}h_0\in\mathrm{GL}_n(\widehat{\mathcal O})$,
and $g\widehat{\mathcal O}^n=h_0\widehat{\mathcal O}^n
=\widehat{N_0}$ exactly.
At $P$ the type residue is that of $h_0$. At any other
$P'\in S_T$, its component is that of $h$; a vector in
$N_0$ belongs to $N_{P'}$, so multiply it by $s\notin P'$
to put it in $N$. The old residue identity and equality
of the prescribed type functionals give the new identity
after cancellation by (AD.8).
The identical argument handles every auxiliary residue vector.
The no-conjugate-pair and disjointness assumptions ensure that
these primes lie outside $\{P,\bar P\}$. Finally all primes
in $Q$ are outside this pair, where $g=h$; the exact rational
frame matrices are unchanged by hypothesis. Thus the frame
field also holds.

Since $h\varepsilon=g$, this proves the complete supported
modification conclusion of \texttt{exists\_splice}, from
the supplied realizer $h_0$. The source obtains $h_0$ from
\texttt{exists\_isPtOf\_cm}; that geometric supplier is
kept explicit. For an actual up/down move one must still
prove its stated equality of local lattices and decorations.
The support conclusion alone does not prove those premises.

\paragraph{AD4: install the exact rational frames at finitely many primes.}
First justify the source's lattice-recovery step. For any
$x\in\mathrm{GL}_n(\mathbb A)$ put
$M_P=x_P\mathcal O_{E_P}^n$. Both $x$ and $x^{-1}$ are integral
outside a finite set. Choose a nonzero rational integer $d$ clearing
their entries at the exceptional primes; then, at every $P$,
\[
 d\mathcal O_{E_P}^n\subseteq M_P\subseteq d^{-1}\mathcal O_{E_P}^n.
                                                               \tag{AD.13}
\]
The rational modules $B_P=M_P\cap E^n$ are local at $P$ in the
sense of LP2 and equal $\mathcal O_{E,(P)}^n$ outside that finite
set. They satisfy the bounds required by LP3. Thus
$N=\bigcap_P B_P$ is a finite full lattice with
$N_{(P)}=B_P$ and $d\mathcal O_E^n\subseteq N\subseteq
d^{-1}\mathcal O_E^n$.
For each fixed $P$, use density of $E$ in its \emph{single}
completion $E_P$ to choose rational columns $v_j$ sufficiently
close to $x_Pe_j$ that
\[
             x_P^{-1}(v_1,\ldots,v_n)
                    \in 1+\mathfrak m_P M_n(\mathcal O_{E_P})
                    \subset\mathrm{GL}_n(\mathcal O_{E_P}).
                                                               \tag{AD.14}
\]
The last inclusion follows by reducing the determinant modulo
$\mathfrak m_P$; it also exhibits the needed openness.
The $v_j$ belong to $B_P$ and form an $\mathcal O_{E_P}$-basis
of $M_P$, so the completion of $N_{(P)}$ is exactly $M_P$.
For a finite lattice between $d\mathcal O_E^n$ and
$d^{-1}\mathcal O_E^n$, its integral adelic span is the product
of these completed local modules: multiply by $d$ and apply the
finite Chinese remainder theorem to the submodule modulo $d^2$.
Therefore $\widehat N=x\widehat{\mathcal O}^n$.
The intersection defining $N$ is precisely
$\operatorname{globLat}(x)$, so this proves the literal
\texttt{Carries} assertion used by \texttt{carries\_globLat}.
The approximation was separate at each single prime, not a
rational strong-approximation assertion. For $n=0$ the statement
is immediate.

Let $x$ be unitary and carry the standard lattice to $N$.
Let $Q$ be a finite set containing no conjugate pair.
At every $P\in Q$ supply a rational matrix $M_P\in
\mathrm{GL}_n(E)$ which is a local basis of $N$ there.
Define a new adelic matrix by
\[
 g_P=M_P\quad(P\in Q),\qquad
 g_{\bar P}=\overline{M_P}^{\,-\mathsf T}\quad(P\in Q),
 \qquad g_w=x_w\quad(w\notin Q\cup\bar Q).           \tag{AD.11}
\]
At the conjugate prime the field matrix is transported to that
completion. Formula (AD.1) constructs each entry, so $g$ is
an actual finite adele. The inverse is obtained by the same
formulas with each selected matrix inverted. The two selected
components are conjugate inverse transposes; hence $g$ is
unitary, including the unchanged components.

Put $b=x^{-1}g$, also unitary. At $P\in Q$, $x_P$ and
$M_P$ are bases of the same local lattice, so both $b_P$
and $b_P^{-1}$ are integral by (AD.7). Unitarity gives
\[
            b_{\bar P}=\overline{b_P^{-1}}^{\,\mathsf T}.
                                                               \tag{AD.12}
\]
It follows that $b_{\bar P}$ and its inverse are integral
as well. Away from $Q\cup\bar Q$ they are identity.
Therefore $b\in\mathrm{GL}_n(\widehat{\mathcal O})$ and
$g$ carries the very same lattice $N$. Its selected component
is exactly $M_P$, rather than merely congruent to it.
Outside $Q\cup\bar Q$, $g^{-1}=x^{-1}$, so all pullback
coordinate functions remain identical. This proves
\texttt{exists\_modify} with all its equality clauses.

For the actual decorated-point application, the type and auxiliary
primes avoid $Q\cup\bar Q$. Their residues are consequently
preserved, while the specified rational frames are installed
exactly. Thus the frame portion of the general realization
theorem is now explicit. Its remaining input is a unitary
adele already realizing the self-dual lattice and the finite
type/auxiliary reductions. The source supplies it using
\texttt{unimodGenus}, \texttt{decoPtSurj\_of} and
\texttt{reduction\_surjective\_bis}; those arithmetic
realization arguments are not inferred from (AD.11).
